% ═══════════════════════════════════════════════════════════════════
%  Chapitre 3 — Étude technique
% ═══════════════════════════════════════════════════════════════════
\chapter{Étude technique}

\section*{Introduction}
\addcontentsline{toc}{section}{Introduction}

Ce chapitre expose les choix techniques du projet. Il présente les langages,
les cadriciels et le système de gestion de base de données retenus, ainsi que
les environnements matériel et logiciel de travail. Il décrit ensuite
l'architecture technique de la solution : son découpage en couches, son
architecture de déploiement, et la chaîne de sécurité appliquée à chaque
requête.

Les choix sont ici justifiés plutôt qu'énumérés. Une technologie ne vaut que
par le problème qu'elle résout dans un contexte donné, et plusieurs des options
retenues n'auraient pas été les mêmes pour un projet d'une autre nature.

% ─────────────────────────────────────────────────────────────────
\section{Technologies de développement}

\subsection{Langages}

\paragraph{TypeScript.} L'ensemble du code applicatif --- API, back-office et
application mobile --- est écrit en TypeScript, sur-ensemble typé de
JavaScript développé par Microsoft. Ce choix unique pour les trois applications
a une conséquence directe et précieuse : les types décrivant les données
métier sont écrits une fois et compris partout. Le type d'un rendez-vous
retourné par l'API est le même côté serveur, côté back-office et côté
application mobile.

Le typage statique joue par ailleurs un rôle de filet de sécurité dans un
projet qui n'a pas de tests d'interface : le back-office et l'application
mobile n'ont pas de suite de tests automatisés, et leur garantie principale est
la compilation. Un champ renommé côté serveur produit une erreur de compilation
côté client, avant tout déploiement.

\paragraph{SQL.} Le langage SQL n'est presque jamais écrit à la main dans le
projet --- Prisma le génère --- à une exception près, précisément là où il est
irremplaçable : la prise du verrou consultatif qui sérialise les réservations
concurrentes, exprimée par un appel direct à \texttt{pg\_advisory\_xact\_lock}.
Aucune abstraction ne fournit cette primitive, et c'est elle qui garantit
qu'un créneau n'est jamais vendu deux fois.

\subsection{Cadriciels et bibliothèques}

\paragraph{NestJS.} L'API repose sur NestJS, cadriciel Node.js structuré autour
de modules, de contrôleurs et de services, et bâti sur l'injection de
dépendances. Sa structure imposée est un atout dans un projet mené seul : elle
évite les choix d'organisation arbitraires et rend chaque module immédiatement
localisable. Ses gardes, filtres, intercepteurs et tubes de validation
permettent en outre de traiter les préoccupations transversales --- droits
d'accès, validation, traduction, gestion des erreurs --- à un seul endroit,
plutôt que répétées dans chaque contrôleur.

\paragraph{Prisma.} Prisma assure la correspondance objet-relationnel. Le
schéma y est déclaré dans un fichier unique, à partir duquel sont générés à la
fois le client typé et les fichiers de migration. Deux propriétés ont pesé dans
ce choix : le client est \emph{typé d'après le schéma}, si bien qu'une colonne
renommée fait échouer la compilation partout où elle était lue ; et les
migrations sont des fichiers SQL versionnés, relus avant application, ce qui
rend le passage en production vérifiable.

\paragraph{React et Vite.} Le back-office est une application web mono-page
écrite en React et servie par Vite. React a été retenu pour la richesse de son
écosystème et sa maîtrise préalable ; Vite pour la rapidité de son serveur de
développement et la simplicité de sa production statique. L'état global ---
réduit à la session authentifiée --- est géré par Zustand, bibliothèque
volontairement minimale, dimensionnée pour ce besoin.

\paragraph{React Native et Expo.} L'application mobile est développée avec
React Native, sous la chaîne d'outils Expo. Ce choix permet de produire une
application Android et iOS à partir d'une base de code unique, tout en
partageant le langage et une partie des conventions avec le back-office. Expo
apporte en outre les briques natives nécessaires --- stockage sécurisé des
jetons, polices, images, dégradés, localisation --- sans avoir à écrire de code
natif.

\paragraph{Bibliothèques notables.} Argon2 pour le hachage des mots de passe,
choisi pour sa résistance aux attaques matérielles ; \texttt{date-fns-tz} pour
les conversions entre heure locale du centre et temps universel ; ExcelJS pour
la génération des exports ; Nodemailer pour l'acheminement des e-mails ; Helmet
pour les en-têtes de sécurité HTTP ; \texttt{class-validator} pour la
validation déclarative des entrées ; i18next pour l'internationalisation de
l'application mobile.

\subsection{Services web}

L'API suit le style d'architecture REST, avec JSON pour format d'échange. Ce
choix s'imposait pour plusieurs raisons convergentes :

\begin{itemize}[leftmargin=1.2cm]
  \item \textbf{L'universalité de HTTP.} Deux clients de natures très
  différentes --- une application native et une application web --- consomment
  la même API sans adaptation.

  \item \textbf{L'absence d'état du protocole.} Chaque requête porte son jeton
  d'authentification ; aucun état de session n'est conservé côté serveur, ce
  qui simplifie considérablement le déploiement et la montée en charge.

  \item \textbf{La légèreté du format.} JSON convient particulièrement à un
  client mobile, dont la bande passante et l'autonomie sont des ressources
  comptées.

  \item \textbf{L'outillage.} La documentation OpenAPI est générée directement
  depuis les décorateurs du code, ce qui garantit qu'elle ne diverge pas de
  l'implémentation.
\end{itemize}

L'API compte plus de quatre-vingts routes réparties en douze contrôleurs. Le
tableau ci-dessous en donne la répartition.

\begin{table}[H]
\centering
\caption{Répartition des routes de l'API par domaine}
\label{tab:routes-api}
\begin{tabularx}{\textwidth}{L{3.1cm} C{1.7cm} X}
\toprule
\entete{Domaine} & \entete{Routes} & \entete{Principales opérations} \\
\midrule
\texttt{/auth} & 9 &
Inscription, connexion, profil courant, vérification d'adresse, renvoi de code,
mot de passe oublié, réinitialisation, rafraîchissement, déconnexion \\

\texttt{/users} & 6 &
Profil, modification, changement de mot de passe, suppression de compte, liste
des utilisateurs, attribution de rôle \\

\texttt{/services} & 10 &
Catalogue public, catégories, détail, gestion complète (administration) \\

\texttt{/products} & 10 &
Catalogue public, catégories, détail, gestion complète (administration) \\

\texttt{/appointments} & 7 &
Disponibilités, réservation, mes rendez-vous, annulation, réservation pour une
cliente, changement de statut, report \\

\texttt{/orders} & 6 &
Commande, mes commandes, détail, annulation, liste (personnel), changement de
statut \\

\texttt{/loyalty} & 23 &
Solde, historique, récompenses, échange, paliers, déblocages, réclamation,
bons, gestion (administration), ajout manuel, journal d'audit \\

\texttt{/opening-\allowbreak hours} & 5 &
Consultation, remplacement des plages, fermetures exceptionnelles \\

\texttt{/settings} & 2 & Consultation et modification des réglages du centre \\
\texttt{/exports} & 3 & Exports Excel des clientes, commandes et rendez-vous \\
\texttt{/uploads} & 1 & Téléversement d'une image de catalogue \\
\texttt{/health} & 1 & Sonde de santé pour la supervision et le déploiement \\
\bottomrule
\end{tabularx}
\end{table}

\subsection{Système de gestion de base de données}

\textbf{PostgreSQL 16} a été retenu comme système de gestion de base de
données. Le choix d'un SGBD relationnel ne faisait pas débat --- les données du
projet sont fortement structurées et fortement liées --- mais celui de
PostgreSQL plutôt que d'une autre base relationnelle mérite d'être justifié.

Trois de ses propriétés sont directement exploitées :

\begin{itemize}[leftmargin=1.2cm]
  \item \textbf{Les verrous consultatifs.} C'est la primitive qui rend correcte
  la réservation concurrente. Elle n'a pas d'équivalent direct et portable
  ailleurs, et elle est utilisée au c\oe ur même du moteur de réservation.

  \item \textbf{Les transactions et l'intégrité référentielle.} Le crédit de
  fidélité, la clôture d'un rendez-vous et le déblocage d'un palier
  s'exécutent dans une transaction unique. Les modes de suppression
  différenciés --- \texttt{Cascade}, \texttt{Restrict}, \texttt{SetNull} selon
  la relation --- portent une partie des règles de conservation de
  l'historique dans le schéma lui-même, où aucun oubli applicatif ne peut les
  contourner.

  \item \textbf{Le type décimal exact.} Les montants sont stockés en
  \texttt{Decimal(10,2)}, et non en nombre flottant. Un montant en virgule
  flottante accumule des erreurs d'arrondi qui finissent par se voir sur un
  total.
\end{itemize}

Le schéma comporte dix-neuf tables et a connu vingt migrations versionnées
depuis l'initialisation du projet.

\subsection{Outils de développement}

\subsubsection{Environnement matériel}

\begin{table}[H]
\centering
\caption{Environnement matériel de développement}
\label{tab:env-materiel}
\begin{tabularx}{\textwidth}{L{4.2cm} X}
\toprule
\entete{Élément} & \entete{Caractéristiques} \\
\midrule
Poste de développement & Ordinateur portable, processeur multic\oe ur, 16 Go
de mémoire vive, disque SSD \\
Système d'exploitation & Windows avec sous-système Linux, et Linux \\
Terminal de test mobile & Téléphone Android physique et émulateur Android \\
Serveur de base de données & Conteneur Docker PostgreSQL 16, exécuté localement \\
\bottomrule
\end{tabularx}
\end{table}

\begin{attention}[Un piège rencontré]
Sous Windows, la plage de ports 2933--3032 est réservée par Hyper-V et par le
sous-système Linux. Le port 3000, choisi par défaut par la plupart des projets
Node.js, y produit une erreur d'accès refusé qui ne désigne pas sa cause. Le
port d'écoute de l'API a donc été rendu configurable par variable
d'environnement --- ce qui s'est révélé nécessaire également au déploiement,
la plupart des hébergeurs imposant le leur.
\end{attention}

\subsubsection{Environnement logiciel}

\begin{table}[H]
\centering
\caption{Environnement logiciel de développement}
\label{tab:env-logiciel}
\begin{tabularx}{\textwidth}{L{4.2cm} X}
\toprule
\entete{Outil} & \entete{Usage} \\
\midrule
Visual Studio Code & Éditeur de code des trois applications \\
Node.js 22 & Environnement d'exécution JavaScript côté serveur et outillage \\
Docker et Docker Compose & Base de données locale, puis conteneurisation de
l'API pour le déploiement \\
Git et GitHub & Gestion de versions et hébergement du dépôt \\
GitHub Actions & Intégration continue : compilation et tests à chaque
modification \\
Prisma Studio & Inspection et correction ponctuelle des données en
développement \\
Swagger UI & Documentation interactive de l'API, générée depuis le code \\
Jest et Supertest & Tests unitaires, tests d'intégration et tests de
concurrence \\
ESLint et Prettier & Analyse statique et mise en forme homogène du code \\
Expo Go & Exécution de l'application mobile sur un téléphone réel sans
compilation native \\
Lucidchart & Réalisation des diagrammes UML \\
\bottomrule
\end{tabularx}
\end{table}

% ─────────────────────────────────────────────────────────────────
\section{Architecture technique}

\subsection{Architecture en couches}

Une vue d'ensemble permet de dégager les couches du système. L'application
étant à la fois web et mobile, l'architecture retenue est une architecture en
quatre niveaux : deux clients distincts, un serveur d'application et un serveur
de données.

\figProjet[0.8]{couches.png}{Les couches de l'application}{couches}

\begin{itemize}[leftmargin=1.2cm]
  \item \textbf{La couche de présentation.} Elle encapsule l'interface
  utilisateur. Elle est double : l'application mobile pour les clientes, le
  back-office web pour l'institut. Elle n'implémente \emph{aucune} règle
  métier : elle affiche ce que le serveur lui transmet et lui soumet ce que
  l'utilisatrice saisit. Les contrôles qu'elle effectue --- un champ
  obligatoire, un format d'adresse --- servent uniquement le confort de saisie
  et sont systématiquement rejoués côté serveur.

  \item \textbf{La couche métier.} C'est le serveur d'application NestJS. Elle
  porte l'intégralité des règles : moteur de créneaux, machines à états,
  calcul des points, contrôle des droits, validation des entrées. Elle expose
  ses services sous forme d'API REST.

  \item \textbf{La couche d'accès aux données.} Assurée par Prisma, elle
  traduit les opérations métier en requêtes SQL, gère les transactions et
  garantit le typage des résultats.

  \item \textbf{La couche de persistance.} La base PostgreSQL. Elle ne se
  contente pas de stocker : elle porte des règles d'intégrité --- unicité,
  clés étrangères, modes de suppression --- et fournit les primitives de
  verrouillage qui rendent la concurrence correcte.
\end{itemize}

\subsection{Architecture de déploiement}

\figProjet[0.9]{architecture.png}{Architecture technique et déploiement de la solution}{architecture}

Le déploiement repose sur Docker Compose et enchaîne trois services dans un
ordre imposé :

\begin{enumerate}[leftmargin=1.2cm]
  \item \textbf{PostgreSQL}, doté d'un contrôle de santé et d'un volume
  persistant. Son port n'est publié que sur l'interface locale.

  \item \textbf{Un service de migration}, à usage unique : il applique les
  migrations en attente, puis s'arrête. Il est délibérément séparé de l'API,
  afin qu'une migration en échec se présente comme une étape ratée assortie de
  son message, et non comme une application qui redémarre en boucle sans
  raison apparente. Il utilise la commande \texttt{migrate deploy}, qui
  applique ce qui manque et refuse de toucher à l'existant --- et jamais
  \texttt{migrate dev}, qui peut réinitialiser la base.

  \item \textbf{L'API}, qui ne démarre qu'une fois les migrations appliquées
  avec succès.
\end{enumerate}

L'image de l'API est construite en trois étapes, afin que l'image livrée ne
contienne que ce qui sert à servir --- ni compilateur, ni cadriciel de test,
ni outillage de développement. Le résultat mesuré est de 178 Mo.

\begin{attention}[Une décision de déploiement]
L'API est publiée sur l'interface locale uniquement. C'est délibéré : un
reverse proxy (Caddy ou nginx) doit se placer devant elle pour assurer le
chiffrement HTTPS. Publier le port directement reviendrait à servir l'API en
clair, jetons d'authentification compris. Le jour où ce proxy est mis en place,
un réglage devient indispensable : la limitation de débit compte par adresse
IP, et derrière un proxy toutes les clientes partagent la sienne. Sans ce
réglage, l'institut entier se retrouve sur un compteur unique et l'API commence
à refuser des requêtes au hasard dès qu'il y a du monde --- un symptôme
déroutant, parfait en test avec une seule personne, et que rien dans les
journaux ne rattache à sa cause.
\end{attention}

\subsection{La chaîne de sécurité d'une requête}

La sécurité n'est pas traitée comme une fonctionnalité, mais comme une chaîne
de filtres que toute requête traverse avant d'atteindre la moindre ligne de
logique métier.

\figProjet[0.85]{securite.png}{Chaîne de traitement de sécurité d'une requête}{securite}

\begin{enumerate}[leftmargin=1.2cm]
  \item \textbf{En-têtes de sécurité.} Helmet retire l'en-tête annonçant la
  technologie du serveur et ajoute une douzaine d'en-têtes appliqués par le
  navigateur : transport strict, anti-\textit{sniffing}, anti-\textit{clickjacking},
  politique de sécurité du contenu.

  \item \textbf{Limitation du débit.} Un compteur par adresse IP plafonne
  l'ensemble de l'API à 120 requêtes par minute. Les routes coûteuses portent
  en plus une limite stricte : dix tentatives par minute pour la connexion et
  l'inscription, cinq pour le renvoi de code, et cinq par tranche de cinq
  minutes pour le mot de passe oublié.

  \item \textbf{Contrôle d'origine.} La politique CORS n'autorise que les
  adresses déclarées du back-office.

  \item \textbf{Authentification.} Un garde vérifie la signature et la validité
  du jeton d'accès, sauf sur les routes explicitement publiques.

  \item \textbf{Autorisation.} Un second garde compare le rôle porté par le
  jeton aux rôles exigés par la route. Ce contrôle est \emph{indépendant} de ce
  que l'interface affiche ou masque : une route d'administration atteinte par
  un autre moyen est refusée par le serveur.

  \item \textbf{Validation des entrées.} Un tube de validation global rejette
  toute propriété non déclarée dans le contrat de la route, plutôt que de
  l'ignorer silencieusement, et convertit les valeurs dans leur type attendu.

  \item \textbf{Traitement uniforme des erreurs.} Un filtre global intercepte
  toute exception avant qu'elle n'atteigne le client. Sans lui, les erreurs de
  base de données partaient brutes, avec les noms de tables et de contraintes.
\end{enumerate}

\paragraph{Le stockage des secrets.} Les mots de passe sont hachés avec Argon2,
lauréat de la compétition internationale de hachage de mots de passe. Les
jetons de rafraîchissement et les codes de vérification ne sont jamais stockés
en clair : seule leur empreinte l'est. Une lecture de la base ne livrerait donc
ni mot de passe, ni session utilisable, ni code de réinitialisation.

\paragraph{La défense contre l'énumération par mesure du temps.} Une
vérification de mot de passe avec Argon2 prend un temps mesurable. Si le
serveur répondait immédiatement pour une adresse inconnue et après le calcul
pour une adresse connue, la seule mesure du délai de réponse permettrait de
déterminer quelles adresses possèdent un compte. La connexion exécute donc une
vérification factice sur une valeur aléatoire lorsque l'adresse est inconnue :
les deux cas répondent au même rythme.

\paragraph{Le contrôle des fichiers téléversés.} Les images du catalogue sont
identifiées par leur \emph{signature binaire}, et non par le type déclaré par
le navigateur ni par l'extension du nom de fichier --- deux valeurs qui
viennent du client et se falsifient en une ligne. Trois formats sont acceptés :
JPEG, PNG et WebP. Le SVG en est délibérément absent : c'est du XML susceptible
de contenir du JavaScript, qui s'exécuterait avec les droits du domaine servant
le fichier --- une faille de type XSS ouverte par un simple téléversement.

\paragraph{La validation de la configuration au démarrage.} L'API refuse de
démarrer si sa configuration est incohérente. Le secret de signature des jetons
doit compter au moins trente-deux caractères ; un secret plus court se casse
hors ligne et permet de forger un jeton d'administration sans jamais toucher à
l'API. De même, l'envoi réel d'e-mails est exigé en production : un
environnement de production configuré en mode console n'enverrait aucun code de
vérification, sans qu'aucune erreur ne le signale.

\section*{Conclusion}
\addcontentsline{toc}{section}{Conclusion}

Ce chapitre a présenté les choix techniques du projet et leur justification :
un langage unique pour les trois applications, un cadriciel structurant pour
l'API, un ORM typé adossé à PostgreSQL, et une architecture en quatre niveaux
où la couche métier est seule dépositaire des règles. La chaîne de sécurité
appliquée à chaque requête a été détaillée, ainsi que l'architecture de
déploiement. Le chapitre suivant présente la mise en \oe uvre effective de ces
choix, les tests qui la valident et les interfaces réalisées.
